\section{The exact integrated trace deficit}\label{app:trace}

The matrix comparison needed in the action argument has an exact
remainder.  We give the proof from \cite[Section~3]{MLT}, retaining
both nonnegative terms.  Together they control the commutator and
the algebraic equality relation used in the action argument.
Throughout this appendix, $0<\sigma<1$, $\beta=\sigma-\tfrac12$, and
$r=1/\sigma$, $p=1+\sigma$, $q=1+r$.
Define $c_\sigma>0$ by
\begin{equation}\label{trace:normalization}
 c_\sigma^{-1}
 =\int_0^\infty t^\beta
       \bigl(t^{-1/2}-(t+1)^{-1/2}\bigr)\,dt.
\end{equation}
The integrand is $O(t^{\sigma-1})$ at zero and
$O(t^{\sigma-2})$ at infinity.  Scaling therefore gives
\begin{equation}\label{trace:scaling}
 c_\sigma\int_0^\infty t^\beta
       \bigl(t^{-1/2}-(t+y)^{-1/2}\bigr)\,dt=y^\sigma,
 \qquad y\geq0.
\end{equation}

\begin{theorem}[Exact trace deficit]\label{trace:deficit}
Let $n\geq1$ and let $B,V\in\C^{n\times n}$ be Hermitian.
Suppose that
\begin{equation}\label{trace:symbol}
 Y_s=(sI-B)^2,\qquad X_s=Y_s+V\geq bI
 \quad(s\in\R)
\end{equation}
for some $b>0$.  Set
\begin{equation}\label{trace:resolvents}
 R(t)=\frac1\pi\int_\R(X_s+tI)^{-1}\,ds,
 \qquad P(t)=t^{-1/2}I-R(t),\qquad t>0,
\end{equation}
and
\begin{equation}\label{trace:N}
 N=c_\sigma\int_0^\infty t^\beta P(t)\,dt.
\end{equation}
These integrals converge absolutely in matrix norm, and $R(t)>0$.
There are continuous Hermitian matrices $L(t)=R(t)^{-1}>0$ and
$Z(t)$ satisfying
\begin{equation}\label{trace:LZ}
 V+tI=L^2+Z^2+i[B,Z],\qquad
 LZ+ZL+i[B,L]=0.
\end{equation}
If $N\geq0$, then
\begin{equation}\label{trace:deficit-identity}
 \begin{split}
 \tr(NV)-\tr(N^q)
 =c_\sigma\int_0^\infty t^\beta\Bigl(
    \tr(RZ^2)
    +q\bigl\|R^{-1/2}-R^{1/2}(tI+N^r)^{1/2}\bigr\|_{\HS}^2
    \Bigr)\,dt.
 \end{split}
\end{equation}
In particular, $\tr(NV)\geq\tr(N^q)$.
\end{theorem}

The hypothesis concerns the full symbol $X_s$; the perturbation
$V$ itself need not be positive.  The proof first recovers $L$ and
$Z$ from solutions on the two half-lines, then evaluates a logarithmic
integral, and finally completes a matrix square.

\subsection{Half-line boundary maps}

Quadratic growth in $s$, together with \eqref{trace:symbol}, gives
an integrable bound for $X_s^{-1}$.  Hence $R(t)>0$ and
$R(t)=O(1)$ as $t\downarrow0$.  Diagonalization of $B$ gives
\[
 \frac1\pi\int_\R(Y_s+tI)^{-1}\,ds=t^{-1/2}I.
\]
For $t\geq1$, both $X_s+tI$ and $Y_s+tI$ dominate
$c(s^2+t)I$, with $c>0$ independent of $s,t$.
The resolvent identity bounds their inverse difference by
$C(s^2+t)^{-2}$, and therefore
\begin{equation}\label{trace:P-bounds}
 \|P(t)\|_{\op}=
 \begin{cases}
  O(t^{-1/2}),&t\downarrow0,\\
  O(t^{-3/2}),&t\to\infty.
 \end{cases}
\end{equation}
These bounds prove the asserted convergence of $N$.

Fix $t>0$ and put $\nabla=\partial_x-iB$.  We shall use the
equation
\begin{equation}\label{trace:ode}
 (-\nabla^2+V+tI)\psi=0.
\end{equation}
Matrix Riccati equations also appear in the commutation method for
Schr\"odinger inequalities \cite{BenguriaLoss2000}.  Here their role
is to compute the integrated resolvent through boundary data.

For each $v\in\C^n$ there is a unique $H^1$ solution of
\eqref{trace:ode} on either half-line with $\psi(0)=v$.
Indeed, the sesquilinear form
\[
 a_t(\phi,\psi)=\int
  \bigl(\langle\nabla\phi,\nabla\psi\rangle
       +\langle\phi,(V+tI)\psi\rangle\bigr)\,dx,
\]
with inner products linear in the second variable, is coercive on
$H^1(\R;\C^n)$: its Fourier symbol $X_s+tI$ dominates
$c_t(1+s^2)I$.  Zero extension gives coercivity on the trace-zero
subspace of either half-line.  Correcting any $H^1$ extension of
$v$ by the Riesz representation theorem gives the unique weak
solution.  The equation implies $\psi''\in L^2$, so this solution
lies in $H^2$; both $\psi$ and $\psi'$ vanish at the infinite end.

Define the right and left boundary matrices $S,T$ by
\begin{equation}\label{trace:boundary}
 \nabla\psi(0+)=-Sv,\qquad \nabla\psi(0-)=Tv.
\end{equation}
Integration by parts gives
$a_t(\psi_v,\psi_w)=\langle v,Sw\rangle$ on the right and
$a_t(\psi_v,\psi_w)=\langle v,Tw\rangle$ on the left.  Thus
$S,T$ are Hermitian.  Translation and uniqueness show that
$\nabla\psi=-S\psi$ throughout the right half-line and
$\nabla\psi=T\psi$ throughout the left half-line.
Substituting these relations into \eqref{trace:ode}, and allowing
arbitrary boundary data, gives
\begin{equation}\label{trace:riccati}
 V+tI=S^2+i[B,S]=T^2-i[B,T].
\end{equation}

Glue the two solutions with boundary value $v$ at zero.
The resulting continuous function has derivative jump
$-(S+T)v$, and hence satisfies
\[
 (-\nabla^2+V+tI)\psi=\delta_0(S+T)v.
\]
Fourier inversion, with forward kernel $e^{-isx}$, gives
\[
 v=\frac1{2\pi}\int_\R(X_s+tI)^{-1}\,ds\,(S+T)v
   =\tfrac12 R(t)(S+T)v.
\]
The inverse symbol is $O(s^{-2})$, so its Fourier inverse is
continuous and this evaluation at zero is justified.  Therefore
\[
 L:=\tfrac12(S+T)=R^{-1}>0,\qquad Z:=\tfrac12(S-T)
\]
satisfy \eqref{trace:LZ}, by adding and subtracting
\eqref{trace:riccati}.  The second equation in \eqref{trace:LZ}
determines $Z$ uniquely from $L$, since $L>0$; in an eigenbasis
of $L$, the map $Z\mapsto LZ+ZL$ multiplies its $(j,k)$ entry
by the positive number $L_{jj}+L_{kk}$.  This also proves the
continuity of $Z(t)$, because $R(t)$, and hence $L(t)$, is continuous.

The diagonal entries of that same equation give $Z_{jj}=0$
in an eigenbasis of $L$, so $\tr Z=0$.  Multiplication on both
sides by $R=L^{-1}$ gives
$i[R,B]=-(ZR+RZ)$.  Cyclicity of the trace now yields
\begin{equation}\label{trace:RV}
 \tr\bigl(R(V+tI)\bigr)=\tr L-\tr(RZ^2).
\end{equation}
In particular, the correction $\tr(RZ^2)=\tr(ZRZ)$ is nonnegative.
We will also need the location of the spectrum of $B+iS$.
The right solutions obey $\psi'=i(B+iS)\psi$.
For an eigenvector with eigenvalue $z$, the solution is $e^{izx}v$;
its square integrability forces $\operatorname{Im}z>0$.

\subsection{A logarithmic integral}

The next step converts the boundary matrices into a scalar identity
that can be integrated in $t$.  Define
\begin{equation}\label{trace:F}
 F(t)=\frac{\tr V}{\sqrt t}
       -\frac1\pi\int_\R
          \log\frac{\det(X_s+tI)}{\det(Y_s+tI)}\,ds.
\end{equation}
The determinants are positive, and the logarithmic ratio is
$O(|s|^{-2})$ at infinity.  Its derivative in $t$ is the trace
of the resolvent difference, with an integrable bound locally
uniform for $t>0$.  Thus the integral converges absolutely and
$F$ is continuously differentiable.

The first identity in \eqref{trace:riccati} factors the symbol as
\[
 X_s+tI=(sI-B+iS)(sI-B-iS).
\]
If $u_j+iv_j$ are the eigenvalues of $B+iS$, counted with
algebraic multiplicity, then $v_j>0$ and
\[
 \det(X_s+tI)=\prod_{j=1}^n\bigl((s-u_j)^2+v_j^2\bigr).
\]
The denominator in \eqref{trace:F} has the same form, with centers
the eigenvalues of $B$ and widths $\sqrt t$.
Translations of the individual logarithms change their
symmetric-cutoff integrals by $o(1)$: for
$h(s)=\log(s^2+v^2)$,
\[
 \frac{d}{da}\int_{-M}^M h(s-a)\,ds
       =h(M+a)-h(M-a)\longrightarrow0
\]
uniformly for bounded $a$.  After centering the factors, the identity
\[
 \int_\R\log\frac{s^2+v^2}{s^2+w^2}\,ds=2\pi(v-w),
 \qquad v,w>0,
\]
follows by differentiation in $v$ and integration from $w$ to $v$.
Since $\sum_jv_j=\tr S=\tr L$, we obtain
\begin{equation}\label{trace:F-evaluation}
 F(t)=\frac{\tr V}{\sqrt t}-2\tr(L-\sqrt t I).
\end{equation}

Set
\begin{equation}\label{trace:G}
 G(t)=\tr(L-2\sqrt t I+tR).
\end{equation}
Both functions are nonnegative.  For $G$, apply
$\lambda-2\sqrt t+t/\lambda\geq0$ to the positive eigenvalues
$\lambda$ of $L$.  For $F$, take the trace of \eqref{trace:LZ}
and use \eqref{trace:F-evaluation} to obtain
\[
 F(t)=t^{-1/2}\tr\bigl((L-\sqrt t I)^2+Z^2\bigr)\geq0.
\]
Equations \eqref{trace:RV} and \eqref{trace:F-evaluation} give
the identity
\begin{equation}\label{trace:FG}
 \tr(PV)=F+G+\tr(RZ^2).
\end{equation}
Every term on the right is nonnegative.  The bounds
\eqref{trace:P-bounds} therefore show that $F$, $G$, and
$\tr(RZ^2)$ are integrable against $t^\beta\,dt$, and that
$t^{\beta+1}F(t)$ vanishes at both endpoints.

Differentiating \eqref{trace:F} and using
\eqref{trace:F-evaluation}--\eqref{trace:G} gives
\[
 F'(t)=-\frac{\tr V}{2t^{3/2}}+\tr P(t),
 \qquad tF'(t)=-\tfrac12F(t)-G(t).
\]
Integration by parts, with the endpoint limits just proved, yields
$\int_0^\infty t^\beta G\,dt
 =\sigma\int_0^\infty t^\beta F\,dt$.
Consequently, integration of \eqref{trace:FG} gives
\begin{equation}\label{trace:weighted}
 \tr(NV)=q c_\sigma\int_0^\infty t^\beta G(t)\,dt
       +c_\sigma\int_0^\infty t^\beta\tr(RZ^2)\,dt.
\end{equation}
It remains to express the first integral as $\tr(N^q)/q$ plus
a nonnegative remainder.

\subsection{The noncommutative square}

For a positive semidefinite matrix $D$, let
\[
 F_0(t,D)=\frac{\tr D}{\sqrt t}
       -2\tr\bigl((tI+D)^{1/2}-\sqrt t I\bigr).
\]
For a scalar eigenvalue $d\geq0$, the corresponding summand is
$\int_0^d(t^{-1/2}-(t+y)^{-1/2})\,dy$.
Tonelli's theorem and \eqref{trace:scaling} therefore imply
\begin{equation}\label{trace:F0}
 c_\sigma\int_0^\infty t^\beta F_0(t,D)\,dt
     =\frac1p\tr(D^p).
\end{equation}
Now expand the Hilbert--Schmidt square
\begin{align}
 \mathcal D_D(t)
  &:=\bigl\|R^{-1/2}-R^{1/2}(tI+D)^{1/2}\bigr\|_{\HS}^2
       \notag\\
  &=\tr\bigl(R^{-1}+R(tI+D)-2(tI+D)^{1/2}\bigr)
       \label{trace:square}\\
  &=G(t)-\tr(P(t)D)+F_0(t,D).\notag
\end{align}
No commutation is used: the cross terms cancel adjacent factors
$R^{-1/2}R^{1/2}$, and cyclicity changes the remaining quadratic
term into $\tr(R(tI+D))$.
All three terms in the last line are integrable by
\eqref{trace:P-bounds}, \eqref{trace:FG}, and \eqref{trace:F0}.
Thus
\begin{equation}\label{trace:duality}
 c_\sigma\int_0^\infty t^\beta G(t)\,dt
  =\tr(ND)-\frac1p\tr(D^p)
       +c_\sigma\int_0^\infty t^\beta\mathcal D_D(t)\,dt.
\end{equation}
When $N\geq0$, choose $D=N^r$.  Since $rp=q=1+r$,
the first two terms on the right equal $\tr(N^q)/q$.
Substitution into \eqref{trace:weighted} proves
\eqref{trace:deficit-identity} and completes the proof of
Theorem~\ref{trace:deficit}.

If $B$ and $V$ commute, the symbol hypothesis implies $V\geq bI$.
Simultaneous diagonalization then gives
$R(t)=(tI+V)^{-1/2}$, $N=V^\sigma$, and $Z=0$.
Both remainders vanish, as required by scalar equality.
